% fig21d_cubic — a 3x3x3 window of a cubic-lattice tensor network
% (Orús review genre, Ann. Phys. 349, 117 (2014) / arXiv:1306.2164:
% 3D PEPS and 3D classical partition functions):
%
%   Z = tTr  prod_{r,c,h}  T^{[r,c,h]}_{l r u d f b}
%
% There is one six-leg tensor at each site.  The fixed plane frame carries
% the row and column axes, while the ordered basis carries the receding axis.
% Its integer quarter-pitch step projects to nearly the same length as a row
% step; no page-space length is chosen here.  All 54 nearest-neighbour
% contractions are stated below: 18 along each axis.
\documentclass[border=6pt]{standalone}
\usepackage{tenkz}
\begin{document}
\begin{tenkz}[
    rows={wire,wire,wire}, cols=3, bonds=none,
    west=none, east=none, north=none, south=none,
    frame={plane,basis={
      wire at (0,0), wire at (-3,5), wire at (-6,10)}}]
  % Column-axis contractions (18).  The coloured middle-layer contraction
  % is one of these records, not an additional edge.
  \tnwire{(1,1,1)}{(1,2,1)}
  \tnwire{(1,2,1)}{(1,3,1)}
  \tnwire{(2,1,1)}{(2,2,1)}
  \tnwire{(2,2,1)}{(2,3,1)}
  \tnwire{(3,1,1)}{(3,2,1)}
  \tnwire{(3,2,1)}{(3,3,1)}
  \tnwire{(1,1,2)}{(1,2,2)}
  \tnwire{(1,2,2)}{(1,3,2)}
  \tnwire[species=marked]{(2,1,2)}{(2,2,2)}
  \tnwire{(2,2,2)}{(2,3,2)}
  \tnwire{(3,1,2)}{(3,2,2)}
  \tnwire{(3,2,2)}{(3,3,2)}
  \tnwire{(1,1,3)}{(1,2,3)}
  \tnwire{(1,2,3)}{(1,3,3)}
  \tnwire{(2,1,3)}{(2,2,3)}
  \tnwire{(2,2,3)}{(2,3,3)}
  \tnwire{(3,1,3)}{(3,2,3)}
  \tnwire{(3,2,3)}{(3,3,3)}

  % Row-axis contractions (18).
  \tnwire{(1,1,1)}{(2,1,1)}
  \tnwire{(2,1,1)}{(3,1,1)}
  \tnwire{(1,2,1)}{(2,2,1)}
  \tnwire{(2,2,1)}{(3,2,1)}
  \tnwire{(1,3,1)}{(2,3,1)}
  \tnwire{(2,3,1)}{(3,3,1)}
  \tnwire{(1,1,2)}{(2,1,2)}
  \tnwire{(2,1,2)}{(3,1,2)}
  \tnwire{(1,2,2)}{(2,2,2)}
  \tnwire{(2,2,2)}{(3,2,2)}
  \tnwire{(1,3,2)}{(2,3,2)}
  \tnwire{(2,3,2)}{(3,3,2)}
  \tnwire{(1,1,3)}{(2,1,3)}
  \tnwire{(2,1,3)}{(3,1,3)}
  \tnwire{(1,2,3)}{(2,2,3)}
  \tnwire{(2,2,3)}{(3,2,3)}
  \tnwire{(1,3,3)}{(2,3,3)}
  \tnwire{(2,3,3)}{(3,3,3)}

  % Receding-axis contractions between adjacent basis members (18).
  \tnwire{(1,1,1)}{(1,1,2)}
  \tnwire{(1,1,2)}{(1,1,3)}
  \tnwire{(1,2,1)}{(1,2,2)}
  \tnwire{(1,2,2)}{(1,2,3)}
  \tnwire{(1,3,1)}{(1,3,2)}
  \tnwire{(1,3,2)}{(1,3,3)}
  \tnwire{(2,1,1)}{(2,1,2)}
  \tnwire{(2,1,2)}{(2,1,3)}
  \tnwire{(2,2,1)}{(2,2,2)}
  \tnwire{(2,2,2)}{(2,2,3)}
  \tnwire{(2,3,1)}{(2,3,2)}
  \tnwire{(2,3,2)}{(2,3,3)}
  \tnwire{(3,1,1)}{(3,1,2)}
  \tnwire{(3,1,2)}{(3,1,3)}
  \tnwire{(3,2,1)}{(3,2,2)}
  \tnwire{(3,2,2)}{(3,2,3)}
  \tnwire{(3,3,1)}{(3,3,2)}
  \tnwire{(3,3,2)}{(3,3,3)}
\end{tenkz}
\end{document}
